% ============================================================
% 油气田生产网数据采集系统 — 技术方案
% dgiot_lite 边缘代理 + 协议逆向 + 采集接管
% V1.0 · 2026-07-10
% ============================================================
\documentclass[12pt,a4paper]{ctexart}

% ---- 包 ----
\usepackage[top=2.5cm,bottom=2cm,left=2.5cm,right=2.5cm]{geometry}
\usepackage{graphicx,float}
\usepackage{fancyhdr,lastpage}
\usepackage{hyperref,bookmark}
\usepackage{longtable,booktabs,array}
\usepackage{enumitem}
\usepackage{caption}
\usepackage{xcolor}
\usepackage{tabularx}
\usepackage{tikz}
\usetikzlibrary{positioning,arrows,shapes,fit,calc}

% ---- 字体 ----
\setmainfont{Microsoft YaHei}
\setCJKmainfont{Microsoft YaHei}
\setmonofont{Consolas}

% ---- 页眉页脚 ----
\pagestyle{fancy}
\fancyhf{}
\fancyhead[L]{\small 油气田生产网数据采集系统 — 技术方案}
\fancyhead[R]{\small V1.0}
\fancyfoot[C]{\thepage}
\renewcommand{\headrulewidth}{0.4pt}

% ---- 超链接 ----
\hypersetup{
  colorlinks=true,
  linkcolor=blue,
  urlcolor=cyan,
  pdftitle={油气田生产网数据采集系统技术方案},
  pdfauthor={迪格(杭州)物联科技有限公司}
}

% ---- 自定义 ----
\definecolor{hl}{HTML}{1a5276}
\newcommand{\kpi}[1]{\textbf{\textcolor{hl}{#1}}}
\newcommand{\code}[1]{\texttt{\small #1}}

\begin{document}

% ===== 封面 =====
\begin{titlepage}
\centering
\vspace*{4cm}

{\Huge \textbf{油气田生产网数据采集系统}}\\[12pt]
{\LARGE \textbf{技术方案}}

\vspace{2cm}

{\large 迪格（杭州）物联科技有限公司}\\[6pt]
{\large 2026 年 7 月}

\vspace{1cm}

\begin{tabular}{rl}
\textbf{版本} & V1.0 \\
\textbf{状态} & 调研阶段 \\
\textbf{密级} & 内部 \\
\textbf{日期} & 2026-07-10 \\
\end{tabular}

\vfill
\end{titlepage}

% ===== 目录 =====
\newpage
\tableofcontents
\newpage

% ============================================================
\section{项目背景与现状}
% ============================================================

\subsection{油气田生产网现状}

某油气田生产网络现有架构运行多年，核心为 pSpace 工业自动化平台：

\begin{table}[H]\centering
\caption{生产网核心服务}
\begin{tabular}{p{2.2cm} p{3.2cm} p{1.8cm} p{5.5cm}}
\toprule
\textbf{服务} & \textbf{进程} & \textbf{版本} & \textbf{路径} \\
\midrule
IO 服务器 & IOMonitor.exe & 7.0.1.2 & \texttt{E:\textbackslash pSpaceOnLine\textbackslash} \\
网桥 & CommBridge.exe & 6.3.0.1 & \texttt{E:\textbackslash pSpaceOnLine\textbackslash} \\
\bottomrule
\end{tabular}
\end{table}

网络架构为典型的三层工控模型：

\begin{figure}[H]\centering
\begin{tikzpicture}[
  every node/.style={font=\footnotesize},
  box/.style={draw, rounded corners, minimum width=2.8cm, minimum height=1.3cm, align=center},
  mainbox/.style={box, fill=blue!15, minimum height=3cm, minimum width=5cm},
  arrow/.style={->, >=stealth, thick},
  line/.style={thick},
]

% 主站 (中心)
\node[mainbox] (main) at (0,0) {
  \textbf{11.66.12.131}\\
  \textbf{IO服务器}\\
  IOMonitor 7.0.1\\
  CommBridge 6.3.0\\
  :8889 ↔ :53001 → :2500 ←
};

% IO 服务器
\node[box, fill=green!12] (io) at (-6,0) {
  \textbf{11.66.12.130}\\
  pSpace\\
  \texttt{A11 :8889}
};

% 数据库
\node[box, fill=gray!10] (db) at (0,-4) {
  \textbf{11.66.12.129}\\
  Oracle DB\\
  \texttt{TNS :1521}
};

% 开发机
\node[box, fill=green!8] (dev) at (6,0) {
  \textbf{11.66.191.155}\\
  开发机\\
  dgiot\_lite
};

% DMZ
\node[box, fill=yellow!12] (dmz) at (6,-4) {
  \textbf{11.66.113.78}\\
  DMZ 中转\\
  MQTT :1883
};

% RTU OPC DA 组
\node[box, fill=orange!10] (rtu_opc) at (-6,4.5) {
  \textbf{172.x.x.x}\\
  RTU/PLC\\
  OPC DA (DCOM)\\
  200+ 台
};

% RTU Modbus 组
\node[box, fill=orange!10] (rtu_mod) at (-2,4.5) {
  \textbf{11.248-9.x.x}\\
  RTU Modbus\\
  slave=1,2
};

% 连线
\draw[arrow] (io) -- node[above]{A11 5a5a} (main);
\draw[arrow] (main) -- node[left]{TNS} (db);
\draw[arrow, red, line width=1.2pt] (main) -- node[above,text=red]{UDP 2500 \\ HDLC+TLS 487K} (dev);
\draw[arrow, blue, dashed] (dev) -- node[left]{MQTT推} (dmz);
\draw[arrow] (main) -- node[above left]{DCE/RPC} (rtu_opc);
\draw[arrow] (rtu_opc) -- node[below right]{OPC响应} (main);
\draw[arrow] (main) -- node[right]{Modbus TCP} (rtu_mod);
\draw[arrow] (rtu_mod) -- node[left]{FC3数据} (main);

\end{tikzpicture}
\caption{生产网整体拓扑 — IP + 协议 + 流向}
\end{figure}

\subsection{核心痛点}

\begin{enumerate}[leftmargin=*]
  \item \textbf{IO 服务器为闭源系统}，A11 私有协议无公开文档，无法扩展或替换
  \item \textbf{缺乏协议级可观测性}，无报文分析能力，故障排查依赖厂家
  \item \textbf{数据无法灵活路由}，现有架构紧耦合，难以对接新平台
  \item \textbf{主站与从站强绑定}，无法灰度迁移或并行运行
\end{enumerate}

\subsection{建设目标}

在不影响现有生产的前提下，通过部署边缘代理实现：

\begin{enumerate}[leftmargin=*]
  \item \textbf{协议逆向}：完整解析 A11 私有协议帧结构
  \item \textbf{透明采集}：旁路获取 IO 流量，结构化存储
  \item \textbf{逐步接管}：旁路分析 → 串联代理 → 完全替换，三阶段平滑过渡
  \item \textbf{多协议融合}：A11 + Modbus TCP/RTU + OPC DA 统一采集
\end{enumerate}

% ============================================================
\section{网络环境调研}
% ============================================================

\subsection{IP 地址规划}

\begin{table}[H]\centering
\caption{关键节点}
\begin{tabular}{p{3.5cm} p{4cm} p{4cm} p{3cm}}
\toprule
\textbf{节点} & \textbf{IP} & \textbf{角色} & \textbf{网段} \\
\midrule
IO服务器 (pSpace 管理端) & 11.66.12.131 & IO 管理 + 数据转发 & 11.66.12.0/24 \\
pSpace A11 :8889 & 11.66.12.0/24 \\
Oracle 数据库 & 11.66.12.129 & 关系存储 (:1521) & 11.66.12.0/24 \\
DMZ 中转 & 11.66.113.78 & MQTT + SSH (:59660) & 11.66.113.0/24 \\
开发机 & 11.66.191.155 & 边缘代理部署 & 11.66.191.128/25 \\
RTU/PLC 群 & 172.x.x.x (200+) & 现场设备 & 多子网 \\
\bottomrule
\end{tabular}
\end{table}

\subsection{端口扫描结果}

\begin{table}[H]\centering
\caption{端口矩阵}
\begin{tabular}{p{2.5cm} p{2cm} p{5.5cm} p{4cm}}
\toprule
\textbf{目标} & \textbf{端口} & \textbf{服务} & \textbf{状态} \\
\midrule
11.66.12.130 & 8889 & A11 pSpace & OPEN \\
11.66.12.131 & 135,139 & RPC + NetBIOS & OPEN \\
11.66.12.129 & 1521 & Oracle TNS & OPEN \\
11.66.113.78 & 1883--1983 & MQTT & DMZ 白名单 \\
11.66.113.78 & 59660 & SSH & DMZ 白名单 \\
\bottomrule
\end{tabular}
\end{table}

\subsection{开发机路由配置}

开发机同时接入 WiFi (192.168.3.111) 和以太网 (11.66.191.155)，双默认路由导致企业内网流量走错出口。

\begin{verbatim}
  route add 11.66.0.0 mask 255.255.0.0 11.66.191.251 metric 5 -p
\end{verbatim}

当前路由表：

\begin{table}[H]\centering
\caption{路由表}
\begin{tabular}{p{3cm} p{2.5cm} p{3cm} p{3cm} p{1.2cm}}
\toprule
\textbf{目标} & \textbf{掩码} & \textbf{网关} & \textbf{接口} & \textbf{Metric} \\
\midrule
0.0.0.0 & 0.0.0.0 & 192.168.3.1 (WiFi) & 192.168.3.111 & 20 \\
0.0.0.0 & 0.0.0.0 & 11.66.191.251 & 11.66.191.155 & 291 \\
\textbf{11.66.0.0} & \textbf{255.255.0.0} & \textbf{11.66.191.251} & \textbf{11.66.191.155} & \textbf{5} \\
\bottomrule
\end{tabular}
\end{table}

% ============================================================
\section{协议逆向分析}
% ============================================================

\subsection{抓包概况}

两份抓包提供了互补视角：

\begin{table}[H]\centering
\caption{抓包对比}
\begin{tabular}{p{3cm} p{4cm} p{4cm}}
\toprule
\textbf{项目} & \textbf{7.3.pcapng} & \textbf{7.10.pcapng} \\
\midrule
抓包位置 & 11.66.12.131 (IO服务器) \\
文件大小 & 356.8 MB & 583.74 MB \\
报文总数 & 251,248 & 952,569 \\
TCP:UDP & 93:6 & 47:52 \\
特征 & 纯生产流量基线 & 开发机 (11.66.191.155) 接入后 \\
\midrule
抓包位置 & 11.66.12.131 & 11.66.191.155 \\
文件大小 & 356.8 MB & 583.74 MB \\
报文总数 & 251,248 & 952,569 \\
TCP:UDP & 93:6 & 47:52 \\
\bottomrule
\end{tabular}
\end{table}

\subsection{协议全景}

\begin{table}[H]\centering
\caption{协议分布（7.10 主站侧）}
\begin{tabular}{p{3.5cm} p{2cm} p{2cm} p{5cm}}
\toprule
\textbf{协议} & \textbf{报文数} & \textbf{占比} & \textbf{说明} \\
\midrule
DCE/RPC (OPC DA) & 72,923 & 16.1\% & 端口 135 + 动态端口，200+ RTU \\
A11 (5a5a) & 93,913 & 20.7\% & 端口 8889，含 ACK 短帧 \\
Modbus TCP & 21,255 & 4.7\% & 端口 53001 ↔ RTU \\
Modbus RTU/TCP & 5,233 & 1.2\% & 串口透传 \\
TLS 1.2 & 4,923 & 1.1\% & 加密应用数据 \\
HDLC+TLS (UDP) & 487,432 & 51.2\% & 端口 2500 ↔ 50123 加密隧道 \\
IEC104 & 39 & $<0.1\%$ & 疑似/误判 \\
\bottomrule
\end{tabular}
\end{table}

\begin{table}[H]\centering
\caption{协议分布（7.3 主站视角）}
\begin{tabular}{p{3.5cm} p{2cm} p{2cm} p{5cm}}
\toprule
\textbf{协议} & \textbf{报文数} & \textbf{占比} & \textbf{说明} \\
\midrule
DCE/RPC (OPC DA) & 41,984 & 16.7\% & OPC DA 直连 RTU 设备 \\
A11 (5a5a) & 34,130 & 13.6\% & 主站 ↔ IO 服务器 \\
Modbus TCP & 5,413 & 2.2\% & 2 个从站 (slave 1+2) \\
Oracle TNS & 8,740 & 3.5\% & 数据库查询 \\
TLS 1.2 & 3,324 & 1.3\% & 加密通道 \\
HDLC+TLS (UDP) & 13,137 & 5.2\% & 加密隧道 \\
\bottomrule
\end{tabular}
\end{table}

\subsection{A11 协议帧结构}

\subsubsection{帧格式}

A11 帧以 \code{5a 5a} 魔术字起始，小端序编码：

\begin{verbatim}
  Offset  Size   Field         Description
  ─────────────────────────────────────────
  0       2B     Magic         5a 5a（帧起始标识）
  2       2B     Frame Length  小端序，不含 2B 头部
  4       4B     Flags         控制标志
  8       2B     Message Type  小端序，标识帧类型
  10      N      Payload       载荷数据（可变）
\end{verbatim}

\subsubsection{真实帧实例}

\textbf{实例 1：设备列表查询帧}（1204B, msg\_type=0xF062）

\begin{verbatim}
  通信: 11.66.12.131:62535 → 11.66.12.130:8889    (7.3.pcapng)
  长度: 1204B    来源: 11.66.12.131 (IO服务器))

5a 5a b2 04 01 00 26 00 62 f0 2f 00 09 00 00 0a
00 24 06 00 00 23 00 00 00 5c 43 59 31 43 38 4b
5c 5a 36 31 31 53 59 57 53 5c 5a 44 30 31 30 38
33 38 44 4f 56 58 56 33 30 31 56 43 23 00 00 00
\end{verbatim}

逐字段解码：

\begin{table}[H]\centering
\caption{A11 查询帧字段解码}
\begin{tabular}{p{2cm} p{1.5cm} p{2.5cm} p{6cm}}
\toprule
\textbf{偏移} & \textbf{长度} & \textbf{值} & \textbf{说明} \\
\midrule
0 & 2B & \code{5a 5a} & 帧起始标识 Magic \\
2 & 2B & \code{04B2 = 1202} & 帧长度（不含 2B 头），LE \\
4 & 4B & \code{01 00 26 00} & 控制标志 \\
8 & 2B & \code{0xF062} & 消息类型：设备列表查询 \\
10 & 1194B & ... & 设备路径数据（ASCII 明文） \\
\bottomrule
\end{tabular}
\end{table}

载荷中包含 ASCII 明文的设备路径：

\begin{verbatim}
  \CY1C8K\Z611SYWS\ZD010838DOVXV301VC#
  \CY1C8K\Z611SYWS\ZD010838DOVXV301VO#
  \CY1C8K\Z611SYWS\ZD010838ERR1VVALVEV0V0#
\end{verbatim}

路径层级解析：

\begin{table}[H]\centering
\caption{设备路径层级}
\begin{tabular}{p{3cm} p{3.5cm} p{6cm}}
\toprule
\textbf{层级} & \textbf{示例} & \textbf{含义} \\
\midrule
采油厂/矿 & CY1C8K & 采油一厂八矿 \\
井/站 & Z611SYWS & 某井站编码 \\
设备标签 & ZD010838DOVXV301VC & 含阀门类型 (DOVX)，编号 (V301)，参数 (VC) \\
错误状态 & ERR1VVALVE & 阀门错误状态监测点 \\
\bottomrule
\end{tabular}
\end{table}

\subsubsection{ACK 短帧}

7.10 抓包中发现大量 4--5 字节的 ACK 确认帧（99\% 的 A11 流量）：

\begin{verbatim}
  5a 5a 02 00            ← 4B ACK (flen=2)
  5a 5a 03 00 06         ← 5B ACK+status (flen=3, val=0x06)
\end{verbatim}

短帧分布（flen=23 占 41,279 个，为最常见类型），表明开发机作为数据消费者频繁确认收到的数据。

\subsubsection{消息类型对比}

\begin{table}[H]\centering
\caption{A11 消息类型对比}
\begin{tabular}{p{2.5cm} p{4cm} p{5cm}}
\toprule
\textbf{7.3 (基线)} & \textbf{7.10 (开发机接入)} & \textbf{差异解读} \\
\midrule
0xF062 (设备列表查询) & --- & 主站专属查询 \\
0x3667 (批量数据上报) & --- & IO 服务器 → 主站 \\
0x87B3 (心跳) & --- & IO 服务器 → 主站 \\
0xF050 (主站查询) & --- & 主站专属 \\
--- & 0x0000 (830 帧) & CommBridge 二次封装 \\
--- & jjZZ 内嵌 & A11 子帧封装标记 \\
\bottomrule
\end{tabular}
\end{table}

开发机通过 CommBridge 接收的是二次封装数据（msg\_type=0x0000），内含 jjZZ 子帧，与主站的原始 A11 帧 (msg\_type 多样化) 完全不同。

\subsection{OPC DA (DCE/RPC) 分析}

OPC DA 基于 DCOM，使用 DCE/RPC v5.0 协议：

\begin{table}[H]\centering
\caption{DCE/RPC 操作统计}
\begin{tabular}{p{3cm} p{2cm} p{8cm}}
\toprule
\textbf{操作} & \textbf{次数} & \textbf{说明} \\
\midrule
Bind & 89 & DCOM 对象导出器绑定（全部成功） \\
Request & 24,939 & OPC Read/Write/Subscribe 请求 \\
Response & 16,864 & OPC 数据返回 \\
\bottomrule
\end{tabular}
\end{table}

通信模式：11.66.12.131 通过 DCOM 直连 172.x 网段 RTU 设备，89 个 OPC 连接覆盖 200+ 台现场设备。

OPC UUID 在 DCOM OXID 编码中不可直接搜索，但协议特征（DCE/RPC v5 + 135 端口 + Bind/Request/Response 模式）完全吻合 OPC DA。

\subsection{Modbus TCP 分析}

\begin{table}[H]\centering
\caption{Modbus TCP 操作统计}
\begin{tabular}{p{3cm} p{2cm} p{8cm}}
\toprule
\textbf{项目} & \textbf{值} & \textbf{说明} \\
\midrule
总帧数 & 21,255 & 端口 53001 ↔ RTU 设备 \\
功能码 FC3 & 20,937 & 读保持寄存器（常规轮询） \\
功能码 FC16 & 318 & 写多寄存器（配置/控制下发） \\
从站地址 & slave=1, slave=2 & 2 个 RTU 从站 \\
\bottomrule
\end{tabular}
\end{table}

\textbf{真实帧实例}（12B, FC3 读保持寄存器, 7.3.pcapng）：

\begin{verbatim}
  通信: 11.66.12.131:53001 → 11.249.61.243:41324

  df 05 00 00 00 06 01 03 01 2b 00 04
\end{verbatim}

\textbf{7.10 对比帧}：

\begin{verbatim}
  通信: 11.66.12.131:53001 → 11.248.198.45:7581

  3b 00 00 00 00 06 01 03 03 d4 00 66
\end{verbatim}

\begin{table}[H]\centering
\caption{Modbus TCP 帧字段解码}
\begin{tabular}{p{2cm} p{1.5cm} p{2.5cm} p{6cm}}
\toprule
\textbf{偏移} & \textbf{长度} & \textbf{值} & \textbf{说明} \\
\midrule
0 & 2B & \code{0xDF05} & 事务 ID = 57093 \\
2 & 2B & \code{0x0000} & 协议 ID = Modbus \\
4 & 2B & \code{0x0006 = 6} & 后续字节数 \\
6 & 1B & \code{0x01} & 从站地址 = 1 \\
7 & 1B & \code{0x03} & 功能码：读保持寄存器 \\
8 & 2B & \code{0x012B = 299} & 起始寄存器地址 \\
10 & 2B & \code{0x0004 = 4} & 读取 4 个寄存器 \\
\bottomrule
\end{tabular}
\end{table}

高频寄存器地址段：

\begin{table}[H]\centering
\caption{寄存器地址段}
\begin{tabular}{p{3cm} p{2cm} p{6cm}}
\toprule
\textbf{地址范围} & \textbf{次数} & \textbf{推测参数} \\
\midrule
299--302 & 349 & 模拟量（压力/温度） \\
350--373 & 238 & 批量模拟量 \\
990--993 & 213 & 电参量 \\
4006 & 212 & 状态量 \\
54272--54271 & 55 & 高频采集参数 \\
\bottomrule
\end{tabular}
\end{table}

\subsection{TLS 加密通道}

\begin{table}[H]\centering
\caption{加密流量统计}
\begin{tabular}{p{3cm} p{2cm} p{2cm} p{5cm}}
\toprule
\textbf{通道} & \textbf{报文数} & \textbf{传输层} & \textbf{说明} \\
\midrule
TLS 1.2 (TCP) & 4,923 & TCP & 端口 2212、2500 等 \\
HDLC+TLS (UDP) & 487,432 & UDP & 端口 2500 ↔ 50123 \\
\bottomrule
\end{tabular}
\end{table}

UDP 2500 HDLC+TLS 封装格式：

\begin{verbatim}
  UDP 载荷:
  ┌──┬────────────────────────────────────┬──┐
  │7E│  HDLC 帧 (去 7E 转义)               │7E│
  └──┴────────────────────────────────────┴──┘
              ↓
  ┌──────┬─────────────────────────────────┐
  │ce 56 │  TLS 1.2 Application Data       │
  │ 4B   │  (加密载荷)                      │
  └──────┴─────────────────────────────────┘
\end{verbatim}

该通道是 CommBridge 将 IO 数据转发至开发机的主要路径，51\% 的流量通过此加密隧道传输。

\subsection{设备发现}

\subsubsection{采油厂设备}

A11 载荷中提取到的设备路径：

\begin{itemize}[leftmargin=*]
  \item \code{\textbackslash CY1C8K\textbackslash Z611SYWS} — 采油一厂八矿 611 井站
  \item \code{ZD010838DOVXV301VC/VO} — 阀门开/关状态
  \item \code{ZD010838COMPV301VJ} — 压缩机运行参数
  \item \code{ZD010838ERR1VVALVE} — 阀门故障监测 (V0--V22)
\end{itemize}

对应传感器数值：177.35, 873.44, 214.19, 220.00（浮点数，推测为压力/温度/流量）。

\subsubsection{RTU 设备}

通过 DCE/RPC 和 Modbus TCP 发现：

\begin{table}[H]\centering
\caption{RTU 设备网段分布}
\begin{tabular}{p{4cm} p{2.5cm} p{5cm}}
\toprule
\textbf{子网} & \textbf{设备数} & \textbf{协议} \\
\midrule
172.26.6.0/24 & 60+ & OPC DA (DCOM) \\
172.21.14.0/24 & 40+ & OPC DA \\
172.23.9.0/24 & 30+ & OPC DA \\
172.23.18.0/24 & 30+ & OPC DA \\
172.28.5.0/24 & 40+ & OPC DA \\
11.249.x.x & 2 & Modbus TCP (slave 1+2) \\
\bottomrule
\end{tabular}
\end{table}

% ============================================================
\section{边缘代理方案}
% ============================================================

\subsection{总体架构}

\begin{figure}[H]\centering
\begin{tikzpicture}[
  box/.style={draw, rounded corners, minimum width=3.5cm, minimum height=2.2cm, align=center, font=\small},
  srv/.style={draw, rounded corners, minimum width=5cm, minimum height=3cm, align=left, font=\footnotesize, fill=blue!5},
  arrow/.style={->, >=stealth, thick, font=\footnotesize},
  net/.style={draw, dashed, rounded corners, inner sep=8pt, font=\small}
]
  % 主站
  \node[srv] (main) at (0,0) {
    \textbf{11.66.12.131 (IO服务器))}\\[2pt]
    IOMonitor.exe \texttt{:8889 ↔}\\
    CommBridge.exe\\
    \texttt{UDP :2500 ←}
  };
  % 开发机
  \node[srv] (dev) at (7,0) {
    \textbf{11.66.191.155 开发机}\\[2pt]
    io\_takeover.py \texttt{:18889}\\
    HTTP API \texttt{:19889}\\
    capture\_server \texttt{:8765}\\
    主平台 \texttt{:8000}\\[2pt]
    \textit{SQLite 持久化}
  };
  % IO
  \node[box,fill=green!10] (io) at (0,3) {
    \textbf{11.66.12.130}\\
    pSpace :8889
  };
  % Oracle
  \node[box,fill=gray!10] (oracle) at (0,-3) {
    \textbf{11.66.12.129}\\
    Oracle :1521
  };
  % DMZ
  \node[box,fill=yellow!10] (dmz) at (7,-3) {
    \textbf{11.66.113.78}\\
    DMZ 中转\\
    MQTT :1883-1983
  };

  % 连线
  \draw[arrow] (io) -- node[right]{A11 5a5a \\ 34K--94K帧} (main);
  \draw[arrow] (main) -- node[right]{TNS} (oracle);
  \draw[arrow, red, line width=1.2pt] (main.east) -- node[above,text=red]{UDP HDLC+TLS \\ \textbf{487K帧 (51\%)}} (dev.west);
  \draw[arrow, blue, dashed] (dev.south) -- node[left]{MQTT推送} (dmz.north);

  % RTU 组（外部设备通过 DCE/RPC 和 Modbus 接入）
\end{tikzpicture}
\caption{边缘代理部署架构}
\end{figure}

\subsection{组件清单}

\begin{table}[H]\centering
\caption{边缘代理组件}
\begin{tabular}{p{3.5cm} p{2cm} p{2.5cm} p{4.5cm}}
\toprule
\textbf{组件} & \textbf{端口} & \textbf{语言} & \textbf{功能} \\
\midrule
\code{io\_takeover.py} & 18889 & Python (183行) & A11 TCP 代理 + 帧解析 + 应答 \\
HTTP API & 19889 & Python & 统计/帧/流/DB 查询 \\
\code{capture\_server.py} & 8765 & Python & 多协议抓包分析 \\
主平台 & 8000 & FastAPI & Web 管理后台 \\
SQLite & 本地 & --- & 报文存储 + 时序数据 \\
\bottomrule
\end{tabular}
\end{table}

\subsection{三阶段部署}

\begin{table}[H]\centering
\caption{部署阶段}
\begin{tabular}{p{1.5cm} p{2.5cm} p{3.5cm} p{5cm}}
\toprule
\textbf{阶段} & \textbf{模式} & \textbf{端口} & \textbf{说明} \\
\midrule
Phase 1 & 旁路分析 & 18889（测试） & 接收 CommBridge 转发，被动采集，零影响 \\
Phase 2 & 串联代理 & 18889→转发→130:8889 & 拦截 A11 流量，双向转发，需维护窗口 \\
Phase 3 & 完全接管 & 8889（替换 130） & 自研 A11 协议栈，原 IO 服务器下线 \\
\bottomrule
\end{tabular}
\end{table}

\subsection{测试验证}

\subsubsection{io\_takeover.py 帧收发测试}

\begin{verbatim}
  TX: 5a 5a 08 00 01 00 00 00 50 f0    (查询帧, msg=0xF050)
  RX: 5a 5a 12 00 01 00 00 00 67 36...  (响应帧, msg=0x3667)
\end{verbatim}

\begin{table}[H]\centering
\caption{测试结果}
\begin{tabular}{p{5cm} p{2cm}}
\toprule
\textbf{测试项} & \textbf{结果} \\
\midrule
TCP 监听 :18889 & 通过 \\
A11 帧解析 (5a5a → 0xF050/0x3667) & 通过 \\
帧收发统计 (connections=1, rx=1, tx=1) & 通过 \\
HTTP API :19889 (\code{/api/stats/frames/flows/db}) & 通过 \\
SQLite 持久化 & 通过 \\
多协议识别 (A11/Modbus/IEC104/jjZZ) & 通过 \\
\bottomrule
\end{tabular}
\end{table}

\subsubsection{真 IO 服务器连接}

成功连接 11.66.12.130 (pSpace A11 :8889) A11 :8889为被动模式（仅响应主站查询），重放 pcap 查询帧因会话状态未建立而未能获得响应，需通过交换机镜像或串联模式进一步测试。

\subsection{报文可见性分析}

基于 11.66.12.131 (IO服务器))的 7.10 抓包（952,569 帧），分析边缘代理的报文收发可见性：

\begin{table}[H]\centering
\caption{协议可见性矩阵}
\begin{tabular}{p{2cm} p{2.2cm} p{1.2cm} p{1.5cm} p{1.8cm} p{3.5cm}}
\toprule
\textbf{协议} & \textbf{通信对} & \textbf{帧数} & \textbf{占比} & \textbf{可见性} & \textbf{可解析内容} \\
\midrule
A11 & 130:8889↔131 & 93.9K & 20.7\% & 🟢 完全 & 设备路径、数值、帧类型 \\
Modbus TCP & 131→RTU & 21.3K & 4.7\% & 🟢 完全 & 从站、FC、寄存器值 \\
DCE/RPC & 172.x↔131 & 72.9K & 16.1\% & 🟡 操作码 & Bind/Req/Resp \\
Oracle TNS & 131→129 & 8.7K & 1.9\% & 🟡 包类型 & Connect/Data \\
\midrule
\multicolumn{6}{l}{\textbf{🔴 加密不可见：}} \\
UDP 隧道 & 155↔131:2500 & 487.4K & 51.7\% & 远程桌面 🔴 加密 & 仅 TLS 头 RDP 加密隧道 \\
TCP TLS & 动态 & 4.9K & 1.1\% & 🔴 加密 & 无 \\
\bottomrule
\end{tabular}
\end{table}

\begin{figure}[H]\centering
\begin{tikzpicture}[font=\small]
  % 边框
  \draw[rounded corners] (0,0) rectangle (12,7.5);

  % 第一列: 加密不可见 (54.5%) — 红色
  \fill[red!30] (0.5,0.5) rectangle (4.2,6.8);
  \node at (2.35,6.3) {\textbf{🔴 加密不可见}};
  \node at (2.35,5.7) {54.5\%};
  \node at (2.35,4.8) {\footnotesize UDP 2500 隧道};
  \node at (2.35,4.3) {\footnotesize 487,432 帧};
  \node at (2.35,3.8) {\footnotesize TLS 1.2 加密};

  % 第二列: 协议可识 (18.0%) — 黄色
  \fill[yellow!40] (4.4,0.5) rectangle (7.6,6.8);
  \node at (6.0,6.3) {\textbf{🟡 协议可识}};
  \node at (6.0,5.7) {18.0\%};
  \node at (6.0,4.8) {\footnotesize DCE/RPC OPC DA};
  \node at (6.0,4.3) {\footnotesize 72,923 帧};
  \node at (6.0,3.8) {\footnotesize Oracle TNS};
  \node at (6.0,3.3) {\footnotesize 操作码可见/内容不可见};

  % 第三列: 明文可解析 (25.4%) — 绿色
  \fill[green!40] (7.8,0.5) rectangle (11.5,6.8);
  \node at (9.65,6.3) {\textbf{🟢 明文可解析}};
  \node at (9.65,5.7) {25.4\%};
  \node at (9.65,4.8) {\footnotesize A11 5a5a};
  \node at (9.65,4.3) {\footnotesize 93,913 帧};
  \node at (9.65,3.8) {\footnotesize Modbus TCP};
  \node at (9.65,3.3) {\footnotesize 21,255 帧};
  \node at (9.65,2.8) {\footnotesize 设备路径+寄存器值};

  % 底部总数
  \node at (6.0,0.2) {\textbf{总报文: 952,569 · 边缘代理可直接解析: 仅 25.4\%}};
\end{tikzpicture}
\caption{边缘代理报文可见性}
\end{figure}

\vspace{-8pt}
\textbf{核心结论}：边缘代理可直接使用的 A11 + Modbus 数据仅占 25.4\%。CommBridge 的 UDP 2500 加密隧道承载了 51.7\% 的数据，需通过串联 io\_takeover 或提取 TLS 会话密钥来突破。

\subsection{Modbus TCP 53001 接管方案（计划中）}

200+ RTU 设备全部通过 \code{131:53001} 走 Modbus TCP，覆盖 11.248/11.249/11.250 三个网段。
接管后边缘代理可获取全部 Modbus 寄存器数据。

\begin{table}[H]\centering
\caption{53001 接管三步}
\begin{tabular}{p{1cm} p{4cm} p{3.5cm} p{3.5cm}}
\toprule
\textbf{步骤} & \textbf{操作} & \textbf{前置条件} & \textbf{回退} \\
\midrule
1 & 131 原 IO 端口改为 53002 & 维护窗口 & 改回 53001 \\
2 & 边缘代理监听 53001，转发到 131:53002 & 步骤1完成 & 代理下线 \\
3 & 代理直接采集，原 IO 端口下线 & 测试稳定 & 恢复原端口 \\
\bottomrule
\end{tabular}
\end{table}

\begin{center}
\textbf{当前不做，等测试环境验证后再执行。}
\end{center}

% ============================================================
\section{实施计划}
% ============================================================

\begin{table}[H]\centering
\caption{实施计划}
\begin{tabular}{p{1.2cm} p{2cm} p{3cm} p{3cm} p{3.5cm}}
\toprule
\textbf{周次} & \textbf{阶段} & \textbf{任务} & \textbf{交付物} & \textbf{风险等级} \\
\midrule
第1周 & Phase 1 & 部署旁路代理，积累协议知识 & 报文分析报告 + A11 帧结构文档 & 低 \\
第2周 & Phase 1 & 解析 Modbus/OPC DA 设备数据 & 设备地址映射表 & 低 \\
第3周 & Phase 2 & 串联代理测试，双向转发验证 & 串联测试报告 + 性能基线 & 中 \\
第4周 & Phase 2 & 数据推送到新平台（MQTT→DMZ） & 数据推送验证 & 中 \\
第5周 & Phase 3 & A11 协议栈实现 + 验收测试 & 完整协议栈 + 测试报告 & 高 \\
第6周 & Phase 3 & 完全接管，原 IO 服务器下线 & 交接文档 + 运维手册 & 高 \\
\bottomrule
\end{tabular}
\end{table}

\subsection{Phase 1 部署步骤（旁路模式，15 分钟）}

\begin{enumerate}[leftmargin=*]
  \item 复制 \code{dgiot\_lite.exe} 部署包到开发机
  \item 配置交换机端口镜像 (SPAN) 或利用 CommBridge 现有转发
  \item 启动边缘代理：\code{python io\_takeover.py 18889}
  \item 启动抓包分析：\code{python capture\_server.py}
  \item 验证：\code{curl http://localhost:19889/api/stats}
\end{enumerate}

\subsection{回退方案}

\begin{enumerate}[leftmargin=*]
  \item \textbf{旁路模式}：直接停止代理进程，零影响
  \item \textbf{串联模式}：\code{pkill -f io\_takeover}，恢复原连接
  \item \textbf{接管模式}：停止代理，重启原 IO 服务器（<2 分钟）
\end{enumerate}

% ============================================================
\section{风险与应对}
% ============================================================

\begin{table}[H]\centering
\caption{风险评估}
\begin{tabular}{p{2.5cm} p{1.5cm} p{3cm} p{5cm}}
\toprule
\textbf{风险} & \textbf{等级} & \textbf{影响} & \textbf{缓解措施} \\
\midrule
协议帧误判 & 中 & A11 解析错误 & 旁路验证 7 天后再串联；误判帧入死信队列 \\
TCP 重组遗漏 & 低 & 丢帧 & 缓冲区保护 + 丢帧日志告警 \\
性能不足 & 低 & 采集延迟 & 纯 Python 已满足 288 pkt/s；瓶颈在 I/O 不在 CPU \\
生产中断 & 高 & 停产事故 & 旁路模式零影响；串联一键回退；操作窗口审批制 \\
数据丢失 & 中 & 历史数据缺失 & SQLite WAL 模式 + 定时备份 + 断点续传 \\
加密破解 & 高 & 无法解析 CommBridge 转发 & 旁路镜像 + Modbus/OPC DA 直采做冗余 \\
\bottomrule
\end{tabular}
\end{table}

\subsection{硬刹车规则}

\begin{enumerate}[leftmargin=*]
  \item 所有操作需在审批维护窗口内执行
  \item 生产网只做被动采集，禁止主动写入
  \item 串联代理需提前准备回退脚本并演练
  \item 所有数据本地存储，不上传外网
\end{enumerate}

% ============================================================
\section{系统部署与运行}
% ============================================================

\subsection{WinRM + netsh trace 远程采集管道}

通过 WinRM (Windows 远程管理) 在主站 131 上执行只读抓包，零安装、不影响任何服务。

\begin{verbatim}
  开发机 (155)                    主站 (131)
  live_collector.py  ──WinRM:5985──→ netsh trace start
       │                                  ↓
       │                             网卡抓包 (30s)
       │                                  ↓
       │                             netsh trace stop
       │                                  ↓
       │                             tracerpt 转 CSV
       │←──────── 下载分析 ─────────────┘
\end{verbatim}

\begin{table}[H]\centering
\caption{远程采集管道组件}
\begin{tabular}{p{3cm} p{5cm} p{4.5cm}}
\toprule
\textbf{组件} & \textbf{位置} & \textbf{功能} \\
\midrule
\code{live\_collector.py} & 开发机 & 定时循环抓包 + 解析 + 推仪表盘 \\
WinRM :5985 & 主站 131 & 远程命令执行 (administrator) \\
\code{netsh trace} & 主站 131 & 网卡层抓包，输出 .etl 文件 \\
\code{tracerpt} & 主站 131 & .etl → CSV 转换 \\
\code{capture\_server.py} & 开发机 :8765 & 多协议解析 + HTTP API \\
Vue 仪表盘 & 开发机 :5173 & Web 实时展示 \\
\bottomrule
\end{tabular}
\end{table}

\subsection{OPC DA 测试结论}

\begin{table}[H]\centering
\caption{OPC DA 连接测试汇总}
\begin{tabular}{p{3cm} p{2.5cm} p{6.5cm}}
\toprule
\textbf{方式} & \textbf{结果} & \textbf{原因} \\
\midrule
C\# DCOM & ❌ & COM QueryInterface 全部返回 E\_NOINTERFACE \\
pywin32 DispatchEx & ❌ & IDispatch 不支持 \\
OpenOPC / openopc2 & ❌ & 底层依赖 IDispatch \\
手动注册 COM 接口 & ❌ & DCS 使用厂商私有 OPC 实现 \\
\bottomrule
\end{tabular}
\end{table}

\textbf{结论}：5 套 DCS 系统使用厂商私有非标 COM 接口，不暴露标准 OPC IID。数据通过 IOMan×36 进程在 131 上正常采集。pcap 中 72,923 帧 DCE/RPC 报文即为 OPC DA 数据流，可通过旁路 NDR 解码获取。

\subsection{已部署服务}

\begin{table}[H]\centering
\caption{开发机服务端口}
\begin{tabular}{p{2cm} p{3.5cm} p{6.5cm}}
\toprule
\textbf{端口} & \textbf{服务} & \textbf{说明} \\
\midrule
:5173 & Vue 仪表盘 & 协议饼图 + 报文表 + 字段解析 \\
:8765 & capture\_server & 多协议解析 (A11/Modbus/OPC-DA/IEC104) \\
:18889 & io\_takeover & A11 TCP 代理 (测试模式) \\
:9876 & relay\_server & 报文中继 (定时从主站抓包) \\
\bottomrule
\end{tabular}
\end{table}

% ============================================================
\section{附录}
% ============================================================

\subsection{设备清单}

\begin{table}[H]\centering
\caption{RTU/PLC 设备统计}
\begin{tabular}{p{3cm} p{2cm} p{3.5cm} p{3cm}}
\toprule
\textbf{协议} & \textbf{设备数} & \textbf{网络} & \textbf{说明} \\
\midrule
Modbus TCP & 206 台 & 11.248--250.x (76子网) & slave=1,2; FC3 轮询 \\
OPC DA & 7 端点 & 172.x (5 DCS) & 200+ 子设备 \\
A11 & -- & 130:8889 & pSpace IO 协议 \\
\bottomrule
\end{tabular}
\end{table}

\subsection{生产网操作规则}

\begin{enumerate}[leftmargin=*]
  \item 所有操作通过 WinRM 只读执行，禁止安装任何软件
  \item 抓包仅用 \code{netsh trace}（Windows 内置），零安装零影响
  \item 禁止停止或修改 IoProject / IoMonitor / CommBridge / IOMan / IoCommit 进程
  \item 开发机与主站之间仅单向拉取数据，不推送任何内容
\end{enumerate}

\vspace{1cm}

\begin{center}
\rule{0.5\textwidth}{0.4pt}\\[6pt]
{\small 迪格（杭州）物联科技有限公司 · 2026 年 7 月 · V1.0}
\end{center}

\end{document}
